% Supplementary material for the AAAI working draft.
\documentclass[10pt,letterpaper]{article}
\usepackage[letterpaper,margin=0.78in]{geometry}
\usepackage{times}
\usepackage{amsmath,amssymb,booktabs,graphicx,url}
\title{Supplementary Material for\\From Visual Goal to Denoising Flow:\\
Carrier-Anchored Image CFG for Training-Free Style Transfer}
\author{Anonymous Authors}
\date{}
\begin{document}
\maketitle

\section{Overview}

This supplement gives the complete formal protocol, additional field
conditioning results, and the online image-reward diagnostic. The image-reward
experiment is included as a negative alignment study. It is not an additional
claim of the main method.

\section{Implementation-Exact Objective}

Both native branches receive the same post-scheduler latent $z_t$, timestep,
prompt metadata, and image-token IDs. Let $v_0(z_t)$ be the differentiable
without-reference velocity and $v_s(z_t)$ the detached with-reference velocity.
The score-match implementation first removes the shared state term and computes
\begin{equation}
 \Delta s_{0\rightarrow s}=-\frac{1-\sigma}{\sigma}
 \left(v_0(z_t)-\operatorname{sg}[v_s(z_t)]\right).
\end{equation}
The coefficient-normalized field enters the direct loss. For carrier anchoring,
the goal direction has the opposite sign,
$g_t=-\operatorname{sg}[\Delta s_{0\rightarrow s}]$, and is transported by
the detached content affinity. The base prediction
$\hat x_0^0=z_t-\sigma v_0(z_t)$ remains differentiable, while
\begin{equation}
 x_0^\star=\operatorname{sg}\left[\hat x_0^0+
 \eta\operatorname{RMSNorm}(\mathcal T_A(g_t))\right]
\end{equation}
is a detached pseudo-target. Constructing and detaching two full scores before
their subtraction would leave an unintended identity gradient from the shared
$-z_t/\sigma$ term.

The implementation includes unit checks for scheduler sigma lookup,
velocity--score round trips, zero correction at the degenerate endpoint,
identity transport, row-stochastic affinity, reference-free carrier inputs, and
absence of gradients through teacher targets.

\section{Formal Reproducibility Protocol}

The representative formal pair uses style image `sty/27.png`, prompt `kp0`,
seed 64, and 512 by 512 output resolution. The main and native teacher roots
are `/mnt/local/klein9b` and `/mnt/local/klein9b-base`. We use 25 outer steps,
five nominal inner iterations with the native dynamic schedule, start time 1000,
AdamW with zero latent weight decay, disabled SLT, disabled warm-start AdaIN,
and disabled final AdaIN.

The score-carrier configuration is:

\begin{center}
\begin{tabular}{ll}
\toprule
Argument & Value\\
\midrule
style gradient RMS target & $2.159326726314248\times 10^{-6}$\\
distance weight & 0.3\\
content weight & 0.4\\
carrier temperature & 0.1\\
carrier strength & 0.75\\
carrier high-pass gain & 0\\
carrier loss weight & 0.25\\
optimizer & AdamW, weight decay 0\\
\bottomrule
\end{tabular}
\end{center}

The exact audit command is:

\begin{verbatim}
/home/ubuntu/chengyanli/flux/bin/python3 \\
  trial/online_rl_critic/run_image_reward_grpo.py \\
  --style-id 27 --prompt-id kp0 --seed 64 \\
  --num-steps 25 --iters 5 --start-time 1000 \\
  --style-mode score_carrier \\
  --style-grad-rms-target 2.159326726314248e-06 \\
  --latent-optimizer adamw --latent-weight-decay 0 \\
  --dist-weight 0.3 --content-weight 0.4 \\
  --carrier-temperature 0.1 --carrier-strength 0.75 \\
  --carrier-highpass-gain 0 --carrier-loss-weight 0.25 \\
  --disable-slt --adain-alpha 0 --disable-final-adain \\
  --disable-grpo \\
  --out-prefix audit_formal_score_carrier_s27_kp0_seed64
\end{verbatim}

The unified evaluator uses one manifest for every row and records the CSD,
CLIP-T, DINO, VGG Gram, and layout-leakage checkpoint hashes. The strict audit
obtains CSD 0.689619, matching the historical formal result around 0.688.

\section{Per-Pair Matched Breakdown}

The main paper reports means over three reference--prompt pairs. Every row
below uses seed 64 and the same sampler and loss configuration within a method.

\begin{table}[h]
\centering
\caption{Per-pair comparison used by the main-paper mean.}
\small
\begin{tabular}{llrrrrr}
\toprule
Pair & Method & CSD $\uparrow$ & CLIP-T $\uparrow$ & DINO $\uparrow$ & Gram $\downarrow$ & Leak $\downarrow$\\
\midrule
00/kp1 & K/V & .4695 & .3035 & .0589 & 4.4740 & .3908\\
       & GoalCarrier & .6726 & .3082 & .0755 & 4.6541 & .4048\\
13/kp3 & K/V & .3129 & .3396 & .0438 & 4.9072 & .3290\\
       & GoalCarrier & .3351 & .3344 & .0410 & 5.1414 & .3382\\
27/kp0 & K/V & .4257 & .3428 & .0400 & 2.7084 & .2858\\
       & GoalCarrier & .6908 & .3463 & .0690 & 2.4718 & .3007\\
\bottomrule
\end{tabular}
\end{table}

The result is heterogeneous. Styles 00 and 27 provide most of the CSD gain;
style 13 is nearly neutral and slightly worse on the auxiliary metrics. We
therefore treat this experiment as controlled mechanism evidence, not a broad
benchmark claim. A complete matched multi-seed matrix is the highest-priority
remaining experiment.

\section{Field Conditioning Ablation}

The optional field student uses either a style-ID embedding or a deterministic
reference image code. The image code is formed from the mean and standard
deviation of flattened VAE tokens. The representative field results are:

\begin{table}[h]
\centering
\caption{Reference conditioning field ablation.}
\begin{tabular}{lrrrr}
\toprule
Condition & Scale & MSE & Cosine & Val. MSE\\
\midrule
Style ID & 0.25 & 0.0000976 & 0.4486 & 0.000223\\
Image code & 0.25 & \textbf{0.0000961} & \textbf{0.4988} & 0.000264\\
\midrule
Style ID holdout & 0.10 & 0.0001211 & \textbf{0.3164} & --\\
Image code holdout & 0.25 & \textbf{0.0001206} & 0.3094 & --\\
\bottomrule
\end{tabular}
\end{table}

The image code is a viable real-reference condition, but the holdout result is
not a universal win. The paper therefore treats the code path as a principled
conditioning interface and leaves stronger encoders and larger reference
families as future work.

\section{Online Image Reward Diagnostic}

\subsection{What Was Tested}

The trial freezes a CSD style model and uses it to score candidate VAE decodes.
For one fixed pre-action latent and three update scales
$a\in\{0.5,1.0,1.5\}$, the candidate reward is:

\begin{equation}
 r_t(a)=10\left[\operatorname{CSD}(D(\operatorname{VAE}(z_t+a\Delta z)),x_s)
 -\operatorname{CSD}(D(\operatorname{VAE}(z_t)),x_s)\right]
 +0.1r_{\mathrm{native}}(a).
\end{equation}

The CSD model is frozen and receives no gradient. The score is computed on a
noisy intermediate latent, not on the final sampled image.

\subsection{Formal Results}

All rows below use the exact formal score-carrier protocol above and the same
reference, prompt, seed, evaluator, and initial noise.

\begin{table}[h]
\centering
\caption{Formal online image-reward diagnostic.}
\small
\begin{tabular}{lrrrrr}
\toprule
Method & CSD $\uparrow$ & CLIP-T $\uparrow$ & DINO $\uparrow$ & Gram $\downarrow$ & Leak $\downarrow$\\
\midrule
Baseline & \textbf{0.689619} & 0.337495 & 0.074046 & 2.390120 & \textbf{0.280550}\\
GRPO, every stage & 0.644513 & 0.332246 & 0.066161 & 2.330238 & 0.308398\\
Local reward oracle & 0.681706 & 0.341500 & 0.067942 & \textbf{2.270237} & 0.298368\\
GRPO, late only & 0.683222 & \textbf{0.354927} & \textbf{0.074363} & 2.280977 & 0.304994\\
\bottomrule
\end{tabular}
\end{table}

The late-only schedule activates candidate selection after outer fraction 0.5.
It avoids most early noisy-state damage and improves several auxiliary
metrics, but remains slightly below the formal CSD baseline. The local reward
oracle is also below baseline CSD.

\subsection{Why This Is Not a Best-of-N Contradiction}

The free-lunch argument for best-of-$N$ assumes that all $N$ candidates are
completed final trajectories and that selection uses the same final objective
reported in evaluation. The trial does neither. It performs greedy one-step
selection:

\begin{equation}
 a_t^\star=\arg\max_a r_t(a),
 \qquad
 r_t(a)=\operatorname{CSD}(\text{noisy intermediate decode}).
\end{equation}

The final objective is instead:

\begin{equation}
 J(a_t)=\operatorname{CSD}(\text{final decode after all future steps}).
\end{equation}

In general, $\arg\max_a r_t(a)$ is not $\arg\max_a J(a_t)$. The current
reward can prefer an update that looks more stylistic before later denoising,
but that changes the state in a way that harms the future score-carrier
trajectory. Moreover, the implemented reward includes native shaping and is
therefore not a pure final-CSD oracle.

A true best-of-$N$ supplementary oracle would fork complete remaining
trajectories from one state, finish every branch with the original
score-carrier sampler, decode the final images, and select the final CSD maximum.
That experiment is more expensive but would have the formal free-lunch property.
The present study should therefore be named a \emph{local image-reward oracle}
or \emph{delayed-reward alignment diagnostic}, not best-of-$N$.

\section{Negative Result and Paper Boundary}

The image-reward implementation successfully evaluates the formal score-carrier
candidate loss without changing the teacher or the main source tree. Its
negative CSD result is informative: local noisy-state CSD is not a sufficiently
aligned value function for greedy online action selection. We do not add this
reward to the main method. The main paper claim remains online native
score-carrier refinement, with image-reward analysis reserved for this
supplement.
The protocol audit, manifests, metrics, logs, and isolated runner are under
\texttt{trial/online\_rl\_critic/}. Generated metrics are ignored by git, while
source code and reproducibility documents are committed.

\end{document}
